\documentclass[11pt,a4paper]{article}
\usepackage{forjatex-manual}

\hypersetup{
  colorlinks=true,
  linkcolor=fjblue,
  urlcolor=fjblue,
  pdftitle={Manual do Usuário - ForjaTeX App Suite},
  pdfauthor={Jean Pimenta / Arkhé Code}
}

\begin{document}

\ForjaHeader{../../forjatex-suite/src-tauri/icons/128x128@2x.png}
  {Ecossistema ForjaTeX App Suite $\cdot$ GNU GPL v3.0}
  {Manual do Usuário --- ForjaTeX App Suite}
  {Hub Central, Instalador e Atualizador de Aplicativos \LaTeX{}}
  {Versão 1.2.6}
  {Jean Pimenta $\cdot$ \href{https://forjatex.com.br}{\color{fjsky}forjatex.com.br} $\cdot$ \href{https://github.com/jeanjpimenta/forjatex-suite}{\color{fjsky}GitHub}}

\ForjaTOC{Sumário}

\begin{fjcallout}[Hub Multiplataforma (Windows \& Linux)]
O \textbf{ForjaTeX App Suite} detecta automaticamente o seu sistema operacional e seleciona instaladores nativos \textbf{.MSI / .EXE} no Windows ou pacotes \textbf{.DEB / .AppImage} no Linux.
\end{fjcallout}

\section{O que é o ForjaTeX App Suite}
O \textbf{ForjaTeX App Suite} é o hub central leve do ecossistema ForjaTeX. Ele identifica os aplicativos instalados no seu computador, permite abri-los em um clique e baixa atualizações diretamente do \textbf{GitHub Releases} sem necessidade de assinatura paga ou conta.

\section{Instalando, Abrindo e Atualizando Aplicativos}
\begin{enumerate}
  \item Na grade principal, cada card exibe a versão instalada ou disponível, o ícone do aplicativo, descrição e botão principal (\textbf{ABRIR}, \textbf{INSTALAR} ou \textbf{ATUALIZAR}).
  \item No seletor \textbf{INSTALADOR} no topo, escolha o formato ideal para o seu sistema:
  \begin{itemize}
    \item \textbf{Windows:} \textbf{.MSI (Padrão)} via \code{msiexec /passive} ou \textbf{.EXE (NSIS)}.
    \item \textbf{Linux:} \textbf{.DEB (Ubuntu/Debian)} via \code{pkexec dpkg -i} ou \textbf{.AppImage} portátil salvo em \code{\~/.local/bin}.
  \end{itemize}
  \item Ao clicar em \textbf{INSTALAR} ou \textbf{ATUALIZAR}, a barra de progresso exibe em tempo real os megabytes baixados e a porcentagem concluída.
  \item Use o botão \textbf{...} em cada card para reinstalar manualmente ou consultar as novidades da versão.
\end{enumerate}

\section{Auto-atualização do Próprio Suite}
Quando uma nova versão do \textbf{ForjaTeX App Suite} estiver disponível no manifesto (\code{suite-manifest.json}), uma faixa superior exibirá o botão \textbf{ATUALIZAR E REINICIAR}. Ao clicar, o Suite baixa o novo instalador, aplica a atualização e reabre automaticamente.

\section{Idioma, Tema e Suporte}
\begin{itemize}
  \item \textbf{Idioma:} detectado automaticamente pelo sistema e alternável entre \textbf{PT-BR} e \textbf{EN} no topo da janela.
  \item \textbf{Tema:} alterne entre Claro e Escuro pelo ícone de Lua/Sol.
  \item \textbf{Apoiar \& Pedidos:} clique em \textbf{APOIAR} para acessar Pix (QR Code), PayPal ou enviar sugestões/bugs diretamente para \code{forjatex.com.br}.
\end{itemize}

\section{Glossário}
\begin{description}
  \item[Manifesto JSON (\code{suite-manifest.json})] Arquivo publicado no GitHub Releases que informa as versões mais recentes e URLs dos instaladores Windows e Linux da suíte.
  \item[MSI / NSIS (.exe)] Formatos padrão de instalação no Windows (\code{msiexec} e instalador NSIS por usuário).
  \item[DEB / AppImage] Pacote nativo para Debian/Ubuntu (\code{.deb}) e executável universal portátil para Linux (\code{.AppImage}).
\end{description}

\section{Histórico de Versões (Changelog)}
\begin{itemize}
  \item \textbf{v1.2.6:} Descrições e subtítulos dos aplicativos totalmente traduzidos para o inglês (\code{EN}) e reabertura automática garantida do Suite após concluir a auto-atualização no Windows.
  \item \textbf{v1.2.5:} Melhoria no instalador MSI no Windows: suporte ao código \code{3010} e fallback automático para instalação interativa quando o modo \code{/passive} for bloqueado por SmartScreen/UAC, exibindo o código exato em caso de falha.
  \item \textbf{v1.2.4:} Painel completo de Configurações em 4 abas (Geral, Atualizações, Downloads \& Cache, Sobre/Logs), ordenação e tema persistidos no \code{localStorage}, botão de atualização do Suite no cabeçalho com barra de progresso em MB antes do reinício e link para o Discord oficial.
  \item \textbf{v1.2.3:} Script de auto-atualização corrigido no Windows e verificação automática em segundo plano com botões \textbf{ABRIR} + \textbf{ATUALIZAR} lado a lado.
  \item \textbf{v1.2.1:} Suporte padrão a \code{.msi} (\code{msiexec /passive}) com seletor \code{.MSI / .EXE}, suporte completo a Linux (\code{.DEB / .AppImage}), tradução PT-BR/EN automática e envio direto de feedback para \code{forjatex.com.br}.
  \item \textbf{v1.2.0:} Layout uniforme estilo ArtCraft (área fixa ao maximizar, efeito hover), detecção real de versão instalada, barra de progresso em MB e auto-atualização (\textbf{ATUALIZAR E REINICIAR}).
  \item \textbf{v1.1.0:} Interface em grade uniforme com ícones dedicados e modal de doação Pix/PayPal.
  \item \textbf{v1.0.0:} Lançamento inicial do hub launcher e integrador com GitHub Releases.
\end{itemize}

\end{document}
